\section{Algoritmo Genetico Ibrido (HGA)}

Tra gli algoritmi proposti, la scelta \`e ricaduta sull'\textbf{Algoritmo
Genetico Ibrido (HGA)} \cite{holland1975adaptation, goldberg1989genetic}, ispirato
al framework HGSADC \cite{vidal2012unified}. L'ibridazione con operatori di ricerca
locale (2-opt, Or-opt, Relocate, Exchange) \cite{lin1973effective, or1976traveling}
combina esplorazione globale e sfruttamento locale.

\subsection{Encoding e Algoritmo di Split}

L'algoritmo utilizza un \textbf{encoding a permutazione} dei clienti.
Ogni cromosoma \`e una permutazione $\pi = (\pi_1, \dots, \pi_n)$ decodificata
in un insieme ottimale di route tramite l'\textbf{algoritmo di Split} di
Prins \cite{prins2004simple} (DP con complessit\`a $O(n^2)$):

$$V[i] = \min_{j: \sum_{k=j}^{i} q_{\pi_k} \leq \sigma} \left[ V[j-1] + d_{0,\pi_j} + \sum_{k=j}^{i-1} d_{\pi_k,\pi_{k+1}} + d_{\pi_i,0} \right]$$

\begin{algorithm}[H]
\scriptsize
\caption{Split Algorithm (Prins, 2004)}
\label{alg:split}
\KwIn{Permutazione $\pi$, matrice distanze $d$, domande $q$, capacit\`a $Q$, depot $0$}
\KwOut{Rotte ottimali $\mathcal{R}$ e costo totale}

$cum[0] \gets 0$\;
\For{$i \gets 1$ \KwTo $n-1$}{$cum[i+1] \gets cum[i] + d[\pi_i,\pi_{i+1}]$\;}

$V[0] \gets 0$; $pred[0] \gets -1$\;
\For{$i \gets 1$ \KwTo $n$}{
    $V[i] \gets +\infty$; $load \gets 0$\;
    \For{$j \gets i$ \KwTo $1$}{
        $load \gets load + q[\pi_j]$\;
        \If{$load > Q$}{\textbf{break}\;}
        $rc \gets d[0,\pi_j] + (cum[i] - cum[j]) + d[\pi_i,0]$\;
        \If{$V[j-1] + rc < V[i]$}{$V[i] \gets V[j-1] + rc$; $pred[i] \gets j-1$\;}
    }
}
\Return{$\textsc{Backtrack}(pred,\pi,n),\; V[n]$}\;
\end{algorithm}

\subsection{Popolazione Iniziale e Operatori Genetici}

La popolazione iniziale ($\mu = 81$) combina: Nearest Neighbor Heuristic,
Clarke-Wright Savings \cite{clarke1964scheduling} e permutazioni casuali per
garantire diversit\`a genetica. La \textbf{selezione a torneo} \cite{goldberg1991comparative}
($k=4$) seleziona il migliore tra 4 individui casuali.

Il \textbf{crossover Order Crossover (OX)} \cite{davis1985applying}
($p_c = 0.675$) preserva l'ordine relativo: seleziona due punti di taglio casuali,
copia il segmento centrale da $P_1$, e riempie le posizioni restanti con i geni
di $P_2$ in ordine, saltando i duplicati (Algoritmo~\ref{alg:ox}).

\begin{algorithm}[H]
\scriptsize
\caption{Order Crossover (OX) \cite{davis1985applying}}
\label{alg:ox}
\KwIn{Permutazioni $P_1$, $P_2$ di lunghezza $n$}
\KwOut{Permutazione figlio $C$}

$cx_1 \gets \text{random}(0, n-1)$; $cx_2 \gets \text{random}(0, n-1)$\;
\If{$cx_1 > cx_2$}{$cx_1, cx_2 \gets cx_2, cx_1$\;}

$C \gets$ array di $n$ elementi; $seg \gets \{\}$\;
\For{$i \gets cx_1$ \KwTo $cx_2$}{$C[i] \gets P_1[i]$; $seg.\text{add}(P_1[i])$\;}

$j \gets (cx_2 + 1) \bmod n$\;
\ForEach{$i \in [cx_2+1, \dots, cx_2+n]$}{
    $pos \gets i \bmod n$\;
    \If{$pos \notin [cx_1, cx_2]$}{
        \While{$P_2[j] \in seg$}{$j \gets (j + 1) \bmod n$\;}
        $C[pos] \gets P_2[j]$; $seg.\text{add}(P_2[j])$; $j \gets (j + 1) \bmod n$\;
    }
}
\Return{$C$}\;
\end{algorithm}

La \textbf{mutazione} ($p_m = 0.236$) applica casualmente: Swap (40\%, scambia
due elementi), Insert (30\%, sposta un elemento), Inversion (30\%, inverte un
segmento).

\subsection{Ricerca Locale Granulare}

Quattro operatori sono applicati con probabilit\`a $p_{ls} = 0.259$ in strategia
\textbf{steepest descent}:
\begin{itemize}[nosep]
    \item \textbf{2-opt} e \textbf{Or-opt} (intra-route): eliminano incroci e
    riposizionano segmenti di 1--3 nodi;
    \item \textbf{Relocate} e \textbf{Exchange} (inter-route): spostano/scambiano
    clienti tra route diverse con filtro GLS \cite{toth2003granular} ($\gamma=25$).
\end{itemize}

La Ricerca Locale Granulare limita le mosse inter-route ai soli $\gamma$ vicini
pi\`u prossimi, riducendo la complessit\`a da $O(N^2)$ a $O(N \cdot \gamma)$.
L'Algoritmo~\ref{alg:relocate} illustra il funzionamento dell'operatore Relocate
con filtro GLS in strategia steepest descent.

\begin{algorithm}[H]
\scriptsize
\caption{Relocate con GLS (steepest descent)}
\label{alg:relocate}
\KwIn{Rotte $\mathcal{R}$, distanze $d$, domande $q$, capacit\`a $Q$, intorni $\Gamma$, depot $0$, max iter $m$}
\KwOut{Rotte migliorate $\mathcal{R}^*$}

$loads[r] \gets \sum_{n \in r} q[n] \; \forall r$\;
$best \gets \mathcal{R}$; $bestCost \gets \textsc{Cost}(best)$; $iter \gets 0$\;

\Repeat{$\neg improved \lor (m > 0 \land iter \geq m)$}{
    $iter \gets iter + 1$; $improved \gets \text{false}$;    $bestMoveCost \gets bestCost$\;
    $bestMove \gets \text{null}$\;
    \ForEach{route $r_f$, nodo $v$ in pos. $f$ di $r_f$}{
        $\delta_f \gets d[prev(v),next(v)] - (d[prev(v),v] + d[v,next(v)])$\;
        \ForEach{route $r_t \neq r_f$ con $loads[r_t] + q[v] \leq Q$}{
            \ForEach{pos. $p \in [0, |r_t|]$}{
                Filtro GLS: salta se entrambi i vicini non sono in $\Gamma[v]$\;
                $\delta_t \gets d[pn,v] + d[v,nn] - d[pn,nn]$\;
                \If{$bestCost + \delta_f + \delta_t < bestMoveCost$}{
                    $bestMove \gets (r_f, f, r_t, p)$\;
                }
            }
        }
    }
    \If{$bestMove \neq \text{null}$}{Applica la mossa; aggiorna $loads$ e $bestCost$; $improved \gets \text{true}$\;}
}
\Return{$best$}\;
\end{algorithm}

\subsection{Ciclo Principale}

L'Algoritmo~\ref{alg:hga} mostra il ciclo evolutivo. I migliori $e=4$ individui
sono preservati (elitismo); i restanti sono generati tramite selezione, OX,
mutazione e ricerca locale. Al termine, i duplicati sono rimpiazzati con individui
casuali. L'algoritmo termina dopo $FE = 350{,}000$ valutazioni.

\begingroup
\scriptsize
\begin{algorithm}[H]
\caption{Ciclo Principale dell'HGA}
\label{alg:hga}
\KwIn{Istanza CVRP; parametri: $\mu$, $p_c$, $p_m$, $p_{ls}$, $k$, $e$, $\gamma$, $FE$}
\KwOut{Migliore soluzione $S^*$}

$\mathcal{P} \gets \{\textsc{Split}(\textsc{NN}()), \textsc{Split}(\textsc{CWS}())\}$\;
\While{$|\mathcal{P}| < \mu$}{$\mathcal{P}.\text{append}(\textsc{Split}(\textsc{Random}()))$\;}
$S^* \gets \arg\min_{S \in \mathcal{P}} S.\text{cost}$; $evals \gets |\mathcal{P}|$\;

\While{$evals < FE$}{
    $\mathcal{P}' \gets \text{top-}e(\mathcal{P})$\;
    \While{$|\mathcal{P}'| < \mu$}{
        \If{$\text{random}() < p_c$}{
            $\pi_c \gets \textsc{OX}(\textsc{Tourn}(\mathcal{P},k), \textsc{Tourn}(\mathcal{P},k))$\;
        }\Else{
            $\pi_c \gets \text{flatten}(\textsc{Tourn}(\mathcal{P},k).\text{routes})$\;
        }
        \If{$\text{random}() < p_m$}{$\pi_c \gets \textsc{Mutate}(\pi_c)$\;}
        $S_c \gets \textsc{Split}(\pi_c)$; $S_c \gets \textsc{TwoOpt}(S_c)$\;
        \If{$\text{random}() < p_{ls}$}{
            $S_c \gets \textsc{LocalSearch}(S_c, \gamma)$\;
        }
        $\mathcal{P}'.\text{append}(S_c)$\;
        \If{$S_c.\text{cost} < S^*.\text{cost}$}{$S^* \gets S_c$\;}
    }
    Rimpiazza duplicati in $\mathcal{P}'$\;
    $\mathcal{P} \gets \mathcal{P}'$\;
}
\Return{$S^*$}\;
\end{algorithm}
\endgroup

\subsection{Parametri Riepilogativi}

\begin{table}[H]
\centering
\caption{Parametri dell'HGA --- Configurazione Tuned}
\label{tab:params}
\footnotesize
\begin{tabular}{@{}ll@{}}
\toprule
\textbf{Parametro} & \textbf{Valore} \\
\midrule
Popolazione ($\mu$) & 81 \\
Crossover rate ($p_c$) / Mutation rate ($p_m$) & 0.675 / 0.236 \\
Local search rate ($p_{ls}$) & 0.259 \\
Tournament ($k$) / Elite ($e$) / Granular ($\gamma$) & 4 / 4 / 25 \\
Max valutazioni (FE) & 350{,}000 \\
\bottomrule
\end{tabular}
\end{table}

\subsection{Architettura Software}

L'algoritmo \`e integrato in un'architettura moderna con backend Python
(FastAPI) \cite{fastapi} e frontend React con visualizzazione real-time tramite
WebSocket. Il backend espone endpoint REST per avviare l'ottimizzazione e
recuperare i risultati, mentre il frontend mostra la disposizione geografica
delle route, i grafici di convergenza e le statistiche aggregate. La
parallelizzazione multi-run sfrutta \texttt{ThreadPoolExecutor}, e
l'accelerazione Numba JIT \cite{lam2015numba} gestisce i calcoli critici
(matrici di distanza, Split DP).
